\subsubsection{LCP Array (Kasai)}

\paragraph{Idea intuitiva.}
El arreglo \textbf{LCP} (Longest Common Prefix) almacena la longitud del prefijo comun mas largo entre pares de sufijos consecutivos en el arreglo de sufijos (SA):
\[
\text{LCP}[k] = \text{lcp}(S[sa[k] \dots n-1], \, S[sa[k+1] \dots n-1])
\]
El algoritmo de \textbf{Kasai et al.} calcula este arreglo en tiempo lineal $O(n)$ procesando los sufijos en el orden en que aparecen en la cadena original ($i = 0, 1, \dots, n-1$) en lugar del orden del SA. La observacion fundamental es que al pasar del sufijo que inicia en $i$ al sufijo que inicia en $i+1$, simplemente se elimina el primer caracter; por ende, si el sufijo $i$ compartia $h$ caracteres con su predecesor en el SA, el sufijo $i+1$ compartira al menos $h - 1$ caracteres con su correspondiente predecesor:
\[
h_{i+1} \ge h_i - 1
\]
Esto garantiza que la variable $h$ se decrementa a lo sumo $n$ veces, y por ende solo puede incrementarse a lo sumo $2n$ veces en total, resultando en tiempo $O(n)$.

\paragraph{Propiedades clave (invariantes).}
\begin{itemize}\setlength\itemsep{0pt}
	\item Arreglo inverso \texttt{rank\_[i]}: posicion en el SA del sufijo que inicia en $i$ (\verb|rank_[sa[k]] = k|).
	\item Tamano del arreglo: $n - 1$ elementos para una cadena de longitud $n$.
	\item LCP de dos sufijos arbitrarios $sa[i]$ y $sa[j]$ con $i < j$: corresponde al minimo elemento en el rango:
	\[
	\text{LCP}(sa[i], sa[j]) = \min_{i \le k < j} \text{lcp}[k]
	\]
	calculable en $O(1)$ tras un precomputo RMQ con Sparse Table sobre el arreglo LCP.
\end{itemize}

\paragraph{Codigo clave.}
Construccion del arreglo LCP con el algoritmo de Kasai:
\begin{verbatim}
vector<int> buildLCP(const string& s, const vector<int>& sa) {
    int n = s.size(), h = 0;
    vector<int> rank_(n), lcp(max(0, n - 1), 0);
    for (int i = 0; i < n; ++i) rank_[sa[i]] = i;
    for (int i = 0; i < n; ++i) {
        if (rank_[i] > 0) {
            int j = sa[rank_[i] - 1];
            while (i + h < n && j + h < n && s[i + h] == s[j + h]) h++;
            lcp[rank_[i] - 1] = h;
            if (h > 0) h--;
        }
    }
    return lcp;
}
\end{verbatim}

\paragraph{Complejidad.}
\begin{itemize}\setlength\itemsep{0pt}
	\item \textbf{Tiempo}: $O(n)$ estricto ($h$ se incrementa como maximo $2n$ veces y se decrementa como maximo $n$ veces).
	\item \textbf{Memoria}: $O(n)$ para el arreglo de retorno y el vector auxiliar \verb|rank_|.
\end{itemize}

\paragraph{Donde tocar para modificar.}
\begin{itemize}\setlength\itemsep{0pt}
	\item \textbf{Subcadenas distintas}: la cantidad total de subcadenas unicas en una cadena de longitud $n$ es:
	\[
	\text{distintas} = \frac{n(n + 1)}{2} - \sum_{k=0}^{n-2} \text{lcp}[k]
	\]
	\item \textbf{Subcadena repetida mas larga}: es el valor maximo del arreglo LCP, obtenido con \verb|*max_element(all(lcp))|.
	\item \textbf{Subcadena comun mas larga de dos cadenas}: concatenar $A + \text{'\$'} + B$, calcular SA y LCP, y buscar el maximo $\text{lcp}[k]$ tal que $sa[k]$ y $sa[k+1]$ pertenezcan a cadenas distintas.
	\item \textbf{Consultas LCP entre pares}: montar una Sparse Table sobre \verb|lcp| para responder $\text{LCP}(u, v)$ en $O(1)$.
\end{itemize}

\paragraph{Trampas y errores frecuentes.}
\begin{itemize}\setlength\itemsep{0pt}
	\item \textbf{Sufijo lexicograficamente menor}: el sufijo con \verb|rank_[i] == 0| no tiene predecesor en el SA; se debe omitir con \verb|if (rank_[i] > 0)|.
	\item \textbf{Cadena de longitud 1 o vacia}: para $n \le 1$, el arreglo LCP tiene tamano 0. Usar \verb|max(0, n - 1)| evita accesos negativos.
	\item \textbf{Olvido de \texttt{h--}}: omitir \verb|if (h > 0) h--;| degrada la complejidad a $O(n^2)$.
\end{itemize}

\paragraph{Explicacion linea por linea.}
\textbf{Funcion buildLCP:}
\begin{itemize}\setlength\itemsep{0pt}
	\item \verb|vector<int> rank_(n);| -- vector de permutacion inversa: ubica donde quedo cada sufijo original dentro del SA.
	\item \verb|for(int i=0;i<n;i++) rank_[sa[i]] = i;| -- llena la tabla de rangos: el sufijo que inicia en \verb|sa[i]| esta en la posicion \verb|i|.
	\item \verb|vector<int> lcp(max(0, n - 1), 0);| -- crea el vector para almacenar las $n-1$ adyacencias.
	\item \verb|int h = 0;| -- longitud coincidente inicial garantizada.
	\item \verb|for(int i = 0; i < n; i++)| -- itera sobre los sufijos en orden de aparicion en el texto.
	\item \verb|if(rank_[i] > 0)| -- evalua unicamente si tiene un vecino inmediatamente anterior en el orden lexicografico del SA.
	\item \verb|int j = sa[rank_[i] - 1];| -- \verb|j| es el indice de inicio del sufijo vecino anterior en el SA.
	\item \verb|while(i + h < n && j + h < n && s[i+h] == s[j+h]) h++;| -- extiende la coincidencia desde el desplazamiento $h$ ya garantizado.
	\item \verb|lcp[rank_[i] - 1] = h;| -- guarda el LCP entre dicho vecino y el sufijo actual.
	\item \verb|if(h > 0) h--;| -- invariante de Kasai: el siguiente sufijo $i+1$ compartira al menos $h-1$ caracteres con su predecesor.
\end{itemize}
\textbf{Programa principal:}
\begin{itemize}\setlength\itemsep{0pt}
	\item \verb|auto lcp = buildLCP(s, sa);| -- calcula las longitudes de prefijos comunes adyacentes para ``banana'' (resultado: 0 1 3 0 0 2).
\end{itemize}

\paragraph{Codigo.}
Ver \texttt{cpp.tex}, seccion \textit{LCP Array (Kasai)}.

\vspace{0.5em}
